\scanpage{29}
\mspage{16}

On a un hom. \inserted{$(\deleted{u_{I'_0},u_{I_0}},u_{I'},u_I)$} déduit de $u$ du carré $4_x$ dans le carré $4_y$, et cet homomorphisme lui-même s’insère dans un diagramme, incluant les deux autres déjà signalés de (2). Voici le diagramme pertinent :

\diagram{\begin{tikzcd}[ampersand replacement=\&,column sep=small,row sep=small]
 \&\&X_{I_0}\arrow[rrr,Rightarrow]\arrow[dd,dashed,no head]\&\&\&X_I\arrow[dd]\&\\
 \&X_{I'_0}\arrow[ur]\arrow[rr]\arrow[dd,no head]\&\&X_{I'}\arrow[urr]\arrow[dd,no head]\&\&\&\\
 \&\&Y_{I'_0,I_0}\arrow[rrr,dash pattern=on 2mm off 3mm on 8mm off .2mm,no head]\arrow[dd,very thick,"{\unsure{u^{\wedge}_{I'_0,I_0}|u_{I'_0,I_0}}}"]\&\&\&
 Y_{I'_0,I}\arrow[ddr,very thick,"{u^I_{I'_0,I_0}}"]\arrow[ddd,thick]\&\\
 (5)\&Y_{I'_0}\arrow[ur]\arrow[rr]\arrow[dr,dash dot]\arrow[dddrr,dashed]\&\&
 Y_{I'_0,I'}\arrow[urr,thick]\arrow[ddd]\&\&\&\\
 \&\&Y_{I_0}\arrow[rrrr,thick]\arrow[ddrrrr,dashed]\&\&\&\&Y_{I_0,I}\arrow[dd,thick,"{u_{I_0,I}}"]\\
 \&\&\&\&\&Y_{I',I}\arrow[dr,very thick,"{u_{I',I}}"']\&\\
 \&\&\&Y_{I'}\arrow[urr]\arrow[rrr,dashed]\&\&\&Y_I
\end{tikzcd}}

% slanted
\margin{diagramme \unsure{imbitable|imbittable|indébrouillable|indéchiffrable}, à remplacer par diagramme (6) ci-dessous}

\deleted{\unsure{Dans}} Le cube supérieur est le cube \inserted{de $\unsure{\mathcal M|\mathscr M}$,} à 6 faces cocartésiennes déduit des trois flèches de $X_{I'_0}$ dans $X_{I_0},X_{I'},Y_{I'_0}$. \inserted{On aura besoin de savoir que le \uline{bas} du cube est un carré cocartésien.} La partie antérieure du diagramme est le \unsure{dessin} de \unsure{(2)|(1)} correspondant aux inclusions $I'_0\subset I'\subset I$, la partie derrière \deleted{est} l’analogue pour \unsure{l’inclusion|les inclusions} $I'_0\subset I_0\subset I$. Les deux autres carrés verticaux, en dessous du cube supérieur, sont également cocartésiens, ainsi que le carré oblique en pointillés, lequel n’est autre que le bas $4_y$ du diagramme (4). De ces \unsure{caractères|conditions} du diagramme 5 résulte \deleted{\unsure{que}} aisément (\unsure{quoique|puisque}
% archive sheet number: 28

